% Generated draft: requires mathematical and source review.
\documentclass[11pt]{article}
\usepackage[margin=25mm]{geometry}
\usepackage{fontspec}
\setmainfont{Latin Modern Roman}
\usepackage{amsmath,amssymb,mathtools,unicode-math}
\setmathfont{Latin Modern Math}
\usepackage{enumitem,needspace,xurl,hyperref}
\setlength{\parindent}{0pt}
\setlength{\parskip}{5pt}
\setlength{\emergencystretch}{3em}
\newcommand{\N}{\mathbb{N}}
\newcommand{\Z}{\mathbb{Z}}
\newcommand{\Q}{\mathbb{Q}}
\newcommand{\R}{\mathbb{R}}
\newcommand{\C}{\mathbb{C}}
\newcommand{\F}{\mathbb{F}}
\newcommand{\A}{\mathbb{A}}
\newcommand{\PP}{\mathbb{P}}
\newcommand{\E}{\mathbb{E}}
\newcommand{\eps}{\varepsilon}
\newcommand{\dd}{\,\mathrm{d}}
\begin{document}
\section*{UnsolvedMath ID: 30005199}
\textbf{Source problem code:} catalogue entry
\subsection{Definitions and mathematical statement}
Let \(C,D,E\) be separable \(C^*\)-algebras, let \(D\triangleleft E\) be a closed two-sided ideal, and let \(q:E\to E/D\) be the quotient map. Let \(\Phi,\Psi:C\to E\) be \(*\)-homomorphisms agreeing modulo \(D\), meaning
\[
q\circ\Phi=q\circ\Psi,
\qquad\text{equivalently}\qquad
\Phi(c)-\Psi(c)\in D\quad(c\in C).
\]
Write \([\Phi,\Psi]_{KK}\in KK(C,D)\) for the associated relative Kasparov class, as in the supplied source. Assume that both maps are absorbing relative to \(D\). The supplied excerpt does not define this absorption condition; its precise source-specific meaning must be recovered before this draft is a fully specified statement.

Let \(\mathcal Z\) denote the Jiang--Su algebra. The condition of \(\mathcal Z\)-stability is
\[
E\cong E\otimes_{\min}\mathcal Z.
\]
Under the usual interpretation, asymptotic unitary equivalence through \(\widetilde D\) means that there exists a norm-continuous path
\[
u:[0,\infty)\longrightarrow\mathcal U(\widetilde D)
\]
such that
\[
\lim_{t\to\infty}\bigl\|u(t)\Phi(c)u(t)^*-\Psi(c)\bigr\|=0
\qquad\text{for every }c\in C.
\]
Here \(\mathcal U(\widetilde D)\) denotes the unitary group of the unitization used in the source, acting on the maps through its natural image in an ambient unitization of \(E\). The excerpt does not specify the unitization convention or whether an initial-value condition such as \(u(0)=1\) is intended. The displayed interpretation imposes no initial-value condition and requires confirmation against the full theorem.

\textbf{Question.} For every such \(C,D,E,\Phi,\Psi\), with the source's relative absorption hypothesis retained, does the equivalence
\[
[\Phi,\Psi]_{KK}=0
\quad\Longleftrightarrow\quad
\Phi\text{ and }\Psi\text{ are asymptotically unitarily equivalent through }\widetilde D
\]
remain valid without assuming \(E\cong E\otimes_{\min}\mathcal Z\)? The supplied excerpt asserts this equivalence under \(\mathcal Z\)-stability; the question concerns removing that hypothesis.

The surrounding excerpt also gives a schematic uniqueness assertion for maps \(\phi,\psi\) and a classification proof sketch involving the UCT, adjustments to \(\phi\), reindexations, and intertwining. It explicitly qualifies that sketch with informal omissions. Those maps and their hypotheses are not fully supplied, so that surrounding discussion is not formulated here as an additional theorem.
\subsection{Short English statement}
For separable \(C^*\)-algebras and two relatively absorbing \(*\)-homomorphisms that agree modulo a closed ideal \(D\), ask whether vanishing of their relative \(KK\)-class is equivalent to asymptotic unitary equivalence through \(\widetilde D\), without assuming \(\mathcal Z\)-stability of the ambient algebra. The source's absorption definition and path conventions still need confirmation.
\subsection{Sources}
\begin{itemize}
\item Supplied local UnsolvedMath registry record, permanent ID \(30005199\), title ``Necessity of Z-Stability in KK-Uniqueness,'' Oberwolfach Reports Open Problems collection, exact source problem code \texttt{catalogue entry}. The record gives proposed year 2022 and names no proposer.
\item Supplied report citation: \(C^*\)-Algebras (2023), \emph{Oberwolfach Reports}, pp.\ 2059--2127, DOI 10.4171/owr/2022/36. Original report link: \url{https://doi.org/10.4171/owr/2022/36}. The supplied title year 2023, proposed year 2022, and DOI containing 2022 are retained without resolving their bibliographic relationship.
\item Supplied literature source: ``Tracially amenable actions and purely infinite crossed products,'' \url{https://doi.org/10.1007/s00208-024-02833-9}.
\item Literature assessment dated exactly \texttt{2026-08-22}: the supplied curation labels the problem open and states that subsequent results establish important special cases and exhibit failures of \(\mathcal Z\)-stability in related nonamenable settings, without settling the general necessity question. This is supplied dated curation, not an independently verified assessment; its status and cited claims require independent checking.
\item The registry labels the statement \texttt{corrected\_verified} and reports that it was reconstructed and checked against the official report. These are supplied provenance claims; no source visit or independent verification is asserted in this draft.
\end{itemize}
\end{document}
